\chapter{Curvature gaps and parallel self-shrinkers}\label{sec:shrinker-curvature-gap}
Throughout this section, $X:\Sigma^n\to\mathbb R^{n+1}$ is a smooth two-sided immersion, $n\geq1$, with $A=D\nu$, $H=\operatorname{tr}A$ and $H=\langle X,\nu\rangle/2$. Write $\rho=(4\pi)^{-n/2}e^{-|X|^2/4}$ and $\mathcal D=\Delta-\frac12\nabla_{X^T}$. The local identities follow the mechanism in \cite{CM12}, Lemma 10.8, pp.~808--809. The complete gap result below is proved for hypersurfaces; compare the arbitrary-codimension theorem in \cite{CL12}, Theorem 1.1, whose normalization is $\vec H=-X^\perp$.

\begin{proposition}[Tensor Simons equation on a shrinker]\label{prop:shrinker-tensor-simons}
The shape tensor satisfies
\[
 \mathcal D A=\left(\frac12-|A|^2\right)A,\qquad
 \mathcal D|A|^2=2|\nabla A|^2+(1-2|A|^2)|A|^2.
\]
No completeness, embeddedness or sign condition on $H$ is needed.
\end{proposition}
\begin{proof}
Differentiate $H=\langle X,\nu\rangle/2$ to get $\nabla H=\frac12A(X^T)$. Since $X^T=X-2H\nu$, tangential differentiation gives $\nabla X^T=I-2HA$. In a geodesic frame, Codazzi therefore yields
\[
 \nabla^2H=\frac12\nabla_{X^T}A+\frac12A-HA^2.
\]
The Euclidean case of Theorem \ref{thm:spaceform-simons} reads $\Delta A=\nabla^2H+HA^2-|A|^2A$. Substitution proves the tensor equation. Applying the product rule for the rough Laplacian and drift to $|A|^2$ gives the scalar equation. This also explains the factor $1/2$ associated to the chosen shrinker normalization.
\end{proof}

\begin{lemma}[A weighted curvature Caccioppoli inequality]\label{lem:shrinker-gap-caccioppoli}
Suppose $f=|A|^2\leq1/2$. For every Lipschitz function $\eta$ compactly supported in the interior,
\[
 \int_\Sigma\eta^2|\nabla A|^2\rho
 +\int_\Sigma\eta^2(1-2f)f\rho
 \leq 2\int_\Sigma|\nabla\eta|^2\rho.
\]
Both terms on the left are nonnegative.
\end{lemma}
\begin{proof}
First take a smooth test. Multiply Proposition \ref{prop:shrinker-tensor-simons} by $\eta^2$ and integrate using Lemma \ref{lem:gaussian-integration-parts}. This gives
\[
 2\int\eta^2|\nabla A|^2\rho+\int\eta^2(1-2f)f\rho
 =-2\int\eta\langle\nabla\eta,\nabla f\rangle\rho.
\]
Metric compatibility gives $|\nabla f|\leq2|A||\nabla A|$, including points where $A=0$. Put $K=1/\sqrt2$. Young's inequality bounds the absolute value of the right side by
\[
 \int\eta^2|\nabla A|^2\rho+4K^2\int|\nabla\eta|^2\rho.
\]
Subtract the first term and use $4K^2=2$. Compactly supported Lipschitz tests follow by smooth approximation in coordinate charts on a fixed compact neighborhood. All coefficients there are bounded, so the integrals converge in the required local Sobolev norms. No division by $|A|$ occurs.
\end{proof}

\begin{proposition}[Finite Gaussian area forces parallel curvature at the gap]\label{prop:shrinker-gap-parallel}
Suppose $\Sigma$ is complete and boundaryless, $\int_\Sigma\rho\,d\mu<\infty$, and $|A|^2\leq1/2$. Then $\nabla A=0$ and $(1-2|A|^2)|A|^2=0$ everywhere. On each connected component, $|A|^2$ is identically zero or identically $1/2$.
\end{proposition}
\begin{proof}
Fix a point in a component. Hopf--Rinow makes closed intrinsic balls compact. The distance function supplies compactly supported Lipschitz cutoffs $\eta_R$ equal to one on the radius-$R$ ball, zero outside the radius-$2R$ ball, with $|\nabla\eta_R|\leq2/R$. Lemma \ref{lem:shrinker-gap-caccioppoli} bounds its left side by $8R^{-2}\int_\Sigma\rho$. On any fixed compact subset, both nonnegative integrands consequently have zero integral on letting $R\to\infty$. Smoothness makes both vanish pointwise. Parallelness makes the norm constant on a connected component, and its algebraic equation leaves the two stated values. Finite Gaussian area is the only global integration hypothesis used in this argument.
\end{proof}

\begin{lemma}[The image of a complete parallel hypersurface]\label{lem:complete-parallel-hypersurface}
Let a connected complete boundaryless two-sided hypersurface in Euclidean space have $\nabla A=0$. Its image is either an affine hyperplane or a round cylinder $S^k(r)\times\mathbb R^{n-k}$, $1\leq k\leq n$, up to an ambient rigid motion. This describes the immersed image and does not assert injectivity of the parametrization.
\end{lemma}
\begin{proof}
Parallel transport preserves the symmetric tensor $A$, so its eigenvalues and multiplicities are constant and its eigendistributions are parallel. If $v,w$ are unit vectors in distinct eigendistributions of eigenvalues $\lambda,\mu$, their sectional curvature is zero. Indeed, the curvature operator preserves the eigendistribution of $w$, and hence $\langle R(v,w)w,v\rangle=0$. The Gauss equation identifies this sectional curvature with $\lambda\mu$. Thus distinct nonzero eigenvalues cannot occur.

If all eigenvalues vanish, the normal is constant and $X$ lies in one affine hyperplane. Otherwise let $\lambda\ne0$ be the unique nonzero eigenvalue, of multiplicity $k$. The images of the nullspace under $dX$ form one fixed ambient plane $L$: parallel transport of a null vector along a curve has both zero covariant derivative and zero normal component of its ambient derivative. Let $\Pi$ be orthogonal projection onto $L^\perp$. Differentiation gives
\[
 d\bigl(\Pi X-\nu/\lambda\bigr)(V)=\Pi dX(V)-dX(AV)/\lambda=0.
\]
Therefore $\Pi X-\nu/\lambda=C$ is constant and the image lies in the cylinder $\{z:|\Pi z-C|=1/|\lambda|\}$. Its tangent space agrees with that cylinder's tangent space at every point. In either case $X$ is a local isometry onto an open subset of the indicated connected complete model.

For completeness of the image, join an image point to any point of the model by a finite-length smooth path, subdivided into coordinate neighborhoods. A local isometry lifts it for as long as local inverse charts exist. If a lift stops at a finite parameter, its remaining lengths tend to zero, so it is Cauchy in the complete source and has a limit there. A local inverse chart at that limit continues it, a contradiction. Hence every such path lifts to its endpoint and the entire model is covered. The case $k=n$ has $L=\{0\}$ and gives a sphere.
\end{proof}

\begin{theorem}[The hypersurface curvature gap]\label{thm:shrinker-half-curvature-gap}
Let $X:\Sigma^n\to\mathbb R^{n+1}$ be connected, complete, boundaryless, smooth and two-sided. Suppose it is a self-shrinker, has finite Gaussian area, and satisfies $|A|^2\leq1/2$. Its image is a hyperplane through the origin or
\[
 S^k(\sqrt{2k})\times\mathbb R^{n-k},\qquad 1\leq k\leq n,
\]
after an orthogonal ambient transformation. If $|A|^2<1/2$ at even one point, the image is a hyperplane. Polynomial area growth is a sufficient alternative to the finite-Gaussian-area hypothesis.
Compare the normalization and gap theorem in \cite{CL12}. The proof below treats hypersurfaces and uses finite Gaussian area; it does not import the arbitrary-codimension classification.
\end{theorem}
\begin{proof}
Proposition \ref{prop:shrinker-gap-parallel} and Lemma \ref{lem:complete-parallel-hypersurface} give the plane or cylinder image. A plane has constant normal and $H=0$, so the shrinker equation forces its normal offset from the origin to vanish.

For a cylinder, use the notation of the lemma, with $C\in L^\perp$. Its normal ranges over the full unit sphere of $L^\perp$, and $H=k\lambda$. The shrinker equation becomes
\[
 2k\lambda=\langle C,\nu\rangle+1/\lambda.
\]
The left side and last term are constant. Replacing a normal direction by its opposite gives $C=0$, and then $\lambda^2=1/(2k)$. This is the claimed radius. Every such cylinder has $|A|^2=k\lambda^2=1/2$, ruling it out under strict inequality at a point. Finally Lemma \ref{lem:gaussian-polynomial-moments} implies finite Gaussian area from polynomial area growth. The proof concerns hypersurfaces; it does not silently apply the hypersurface Simons equation in higher codimension.
\end{proof}

\begin{corollary}[Compact and strict versions of the gap]\label{cor:compact-shrinker-curvature-gap}
A connected closed two-sided self-shrinker with $|A|^2\leq1/2$ has image the sphere of radius $\sqrt{2n}$ centered at the origin. In particular every closed such shrinker satisfies $\max|A|^2\geq1/2$, and equality gives that spherical image.
\end{corollary}
\begin{proof}
A closed immersed manifold has finite Gaussian area. Its image is compact, so the plane and the cylinders with a positive-dimensional Euclidean factor in Theorem \ref{thm:shrinker-half-curvature-gap} are impossible. If the maximum were strictly below $1/2$, that theorem would instead force a plane. The sphere has $|A|^2=1/2$ by its principal curvatures.
\end{proof}

\begin{proposition}[A mean-convex weighted curvature estimate]\label{prop:meanconvex-shrinker-weighted-estimate}
On a two-sided shrinker with $H>0$, every smooth $\phi$ compactly supported in the interior satisfies
\[
 \int |A|^2\phi^2\rho\leq\int|\nabla\phi|^2\rho+\frac12\int\phi^2\rho.
\]
This is the curvature estimate underlying \cite{CM12}, Proposition 10.14. It is not the definition of F-stability, which also varies center and scale.
\end{proposition}
\begin{proof}
Tracing Proposition \ref{prop:shrinker-tensor-simons} gives $\mathcal D H=(1/2-|A|^2)H$. The chain rule for the positive function $H$ yields
\[
 \mathcal D\log H=\frac12-|A|^2-|\nabla\log H|^2.
\]
Multiply by $\phi^2$ and use compactly supported weighted integration by parts. Rearranging gives
\[
 \int |A|^2\phi^2\rho=\frac12\int\phi^2\rho
 -\int\phi^2|\nabla\log H|^2\rho
 +2\int\phi\langle\nabla\phi,\nabla\log H\rangle\rho.
\]
Completing the square bounds the last two terms by $\int|\nabla\phi|^2\rho$. Positivity is only needed on a neighborhood of the support of the test; completeness and growth assumptions play no role in this local estimate.
\end{proof}
